\section{Техническое задание}

\subsection{Основание для разработки}

Основанием для разработки является задание на выпускную квалификационную работу бакалавра "Информационно-аналитическая система для РИО вуза.Разработка сервиса "Методические указания"

\subsection{Цель и назначение разработки}

Основной целью работы является разработка и внедрение информационно-аналитической системы для автоматизации процессов проверки, хранения и сопровождения учебно-методических и научных материалов, поступающих в редакционно-издательский отдел вуза.

Информационно-аналитическая система предназначена для преподавателей, являющизся авторами учебно-методических разработок и сотрудников релакционно-издательского вуза.

\subsection{Требования к функциональности}

Информационная система должна обеспечивать выполнение следующих функций:

\begin{enumerate}
	\item {Управление пользователями.} Система должна обеспечивать регистрацию и авторизацию пользователей по логину и паролю. Должно быть реализовано разграничение прав доступа: администратор, редактор, пользователь. Администратор должен иметь возможность управления учётными записями пользователей.
	
	\item {Загрузка документов.} Пользователь должен иметь возможность загружать методические материалы, учебные пособия, монографии и иные документы в формате DOCX через веб-интерфейс системы.
	
	\item {Автоматическая проверка оформления.} После загрузки документа система должна автоматически анализировать файл на соответствие установленным требованиям оформления. Проверка должна включать контроль структуры документа, параметров шрифта, интервалов, отступов, заголовков, таблиц, изображений, списков литературы и иных элементов.
	
	\item {Формирование отчёта об ошибках.} По завершении анализа система должна формировать отчёт, содержащий перечень обнаруженных нарушений, описание ошибок и рекомендации по их устранению.
	
	\item {Создание документа с выделением ошибок.} Система должна создавать копию загруженного документа в формате DOCX, в которой найденные ошибки визуально выделены средствами форматирования.
	
	\item {Хранение корректных документов.} Если загруженный документ соответствует всем требованиям, он должен автоматически сохраняться в централизованной базе данных или файловом архиве системы.
	
	\item {Поиск и просмотр документов.} Система должна обеспечивать просмотр списка ранее загруженных материалов, поиск по названию, автору, дате загрузки, статусу проверки и другим параметрам.
	
	\item {Генерация титульных листов.} Пользователь должен иметь возможность заполнить форму с исходными данными, после чего система должна автоматически сформировать титульный лист по утверждённому шаблону вуза.
	
	\item {Статистика и аналитика.} Система должна формировать сводные данные по количеству загруженных документов, числу успешных проверок, количеству найденных ошибок, активности пользователей и временным показателям обработки материалов.
	
	\item {Логирование действий.} Система должна сохранять сведения о действиях пользователей: загрузка файлов, проверка документов, вход в систему, генерация титульных листов и иные операции.
\end{enumerate}

\subsection{Требования к интерфейсу}

Интерфейс информационной системы представляет собой веб-приложение с единым современным стилем оформления, обеспечивающее удобную и интуитивно понятную навигацию. Пользовательский интерфейс должен быть адаптирован для работы в распространённых браузерах и обеспечивать быстрый доступ к основным функциям системы.

Основные элементы интерфейса включают:

\begin{itemize}
	\item главную страницу системы;
	\item форму авторизации и регистрации пользователей;
	\item личный кабинет пользователя;
	\item страницу загрузки документов на проверку;
	\item страницу просмотра результатов проверки;
	\item страницу архива сохранённых документов;
	\item страницу генерации титульных листов;
	\item административную панель управления пользователями;
	\item модуль просмотра статистики и аналитических данных.
\end{itemize}

На рисунке~\ref{fig:authorization} представлена страница авторизации системы.

\begin{figure}[h]
	\centering
	\includegraphics[width=1\textwidth]{authorization.png}
	\caption{Страница авторизации}
	\label{fig:authorization}
\end{figure}

Как показано на рисунке, форма авторизации содержит поля ввода логина и пароля, кнопку входа в систему, а также элементы перехода к регистрации нового пользователя. После успешной авторизации пользователь получает доступ к функциональным возможностям системы в соответствии со своей ролью.

На рисунке~\ref{fig:profile} представлен внешний вид личного кабинета пользователя.

\begin{figure}[h]
	\centering
	\includegraphics[width=1\textwidth]{profile.png}
	\caption{Личный кабинет пользователя}
	\label{fig:profile}
\end{figure}

Личный кабинет содержит навигационное меню, сведения о текущем пользователе, основные разделы системы, а также быстрый доступ к загрузке документов, архиву файлов и генерации титульных листов.

\clearpage
На рисунке~\ref{fig:filecheck} представлена страница проверки документов.

\begin{figure}[h]
	\centering
	\includegraphics[width=1\textwidth]{filecheck.png}
	\caption{Страница проверки документов}
	\label{fig:filecheck}
\end{figure}

Данная страница содержит форму выбора файла формата DOCX, кнопку запуска автоматической проверки, область отображения текущего статуса обработки и результаты анализа после завершения проверки.

На рисунке~\ref{fig:titlegen} представлена страница генерации титульных листов.

\begin{figure}[h]
	\centering
	\includegraphics[width=1\textwidth]{titlegen.png}
	\caption{Страница генерации титульных листов}
	\label{fig:titlegen}
\end{figure}

Как показано на рисунке, страница содержит форму ввода реквизитов работы: наименование дисциплины, тема работы, сведения об авторе, руководителе и кафедре. После заполнения формы пользователь может сформировать готовый документ в формате DOCX.

На рисунке~\ref{fig:statistic} представлена страница статистики и аналитики.

\begin{figure}[h]
	\centering
	\includegraphics[width=1\textwidth]{statistic.png}
	\caption{Страница статистики}
	\label{fig:statistic}
\end{figure}

Страница статистики содержит сведения о количестве зарегистрированных пользователей, числе загруженных документов, результатах проверок, количестве обнаруженных ошибок, а также иные аналитические показатели деятельности системы.

Интерфейс системы должен обеспечивать логичное расположение элементов управления, единообразие оформления страниц, наглядность представления результатов работы и минимальное количество действий пользователя при выполнении основных операций.


\subsection{Требования к обработке ошибок}

Система должна обеспечивать корректную обработку исключительных ситуаций.

Основные требования:

\begin{itemize}
	\item при возникновении ошибки пользователю должно отображаться понятное сообщение;
	\item внутренние ошибки системы не должны раскрывать служебную информацию;
	\item ошибки должны фиксироваться в журнале событий;
	\item должна быть реализована обработка:
	\begin{itemize}
		\item некорректных входных данных;
		\item ошибок загрузки файлов;
		\item ошибок обработки документов;
		\item ошибок взаимодействия с базой данных.
	\end{itemize}
\end{itemize}

\subsection{Требования к удобству использования}

Система должна обеспечивать удобство и простоту взаимодействия пользователей.

Основные требования:

\begin{itemize}
	\item интерфейс должен быть интуитивно понятным;
	\item количество действий для выполнения основной операции (загрузка и проверка документа) должно быть минимальным;
	\item результаты проверки должны отображаться в структурированном виде;
	\item должна быть реализована визуальная обратная связь (индикаторы загрузки, статусы операций).
\end{itemize}

\subsection{Моделирование вариантов использования}

Для определения функциональных возможностей разрабатываемой системы и взаимодействия пользователей с её компонентами была выполнена модель вариантов использования.

В системе предусмотрены следующие роли пользователей:

\textbf{Администратор} --- обладает полными правами доступа, включая управление пользователями, просмотр статистики, настройку системы, удаление документов и контроль работы пользователей.

\textbf{Редактор} --- осуществляет просмотр загруженных материалов, контроль результатов автоматической проверки, работу с архивом документов и формирование отчётности.

\textbf{Пользователь} --- выполняет загрузку документов, получает результаты проверки, скачивает отчёты, использует модуль генерации титульных листов и просматривает историю своих материалов.

Основные прецеденты (варианты использования) системы:

\begin{enumerate}
	\item Вход в систему.
	\item Регистрация пользователя.
	\item Загрузка документа на проверку.
	\item Автоматическая проверка оформления.
	\item Просмотр отчёта об ошибках.
	\item Скачивание исправленного документа.
	\item Сохранение корректного документа в архив.
	\item Поиск и просмотр ранее загруженных работ.
	\item Генерация титульного листа.
	\item Просмотр статистики работы системы.
	\item Управление пользователями.
	\item Выход из системы.
\end{enumerate}

Для повышения наглядности диаграмма прецедентов была разделена по ролям пользователей. Отдельно представлены варианты использования для пользователя, администратора и редактора.

\clearpage
На рисунке~\ref{fig:diagramUser} представлена диаграмма прецедентов пользователя.

\begin{figure}[h]
	\centering
	\includegraphics[width=1\textwidth]{diagramUser.png}
	\caption{Диаграмма прецедентов пользователя}
	\label{fig:diagramUser}
\end{figure}


На рисунке~\ref{fig:diagramAdmin} представлена диаграмма прецедентов администратора.

\begin{figure}[h]
	\centering
	\includegraphics[width=0.9\textwidth]{diagramAdmin.png}
	\caption{Диаграмма прецедентов администратора}
	\label{fig:diagramAdmin}
\end{figure}

\clearpage
На рисунке~\ref{fig:diagramEditor} представлена диаграмма прецедентов редактора.

\begin{figure}[h]
	\centering
	\includegraphics[width=0.9\textwidth]{diagramEditor.png}
	\caption{Диаграмма прецедентов редактора}
	\label{fig:diagramEditor}
\end{figure}

\subsubsection{Вариант использования «Вход в систему»}

\textbf{Заинтересованные лица и их требования:} пользователь желает получить доступ к функциям системы путём ввода логина и пароля.

\textbf{Предусловие:} пользователь не авторизован; открыта страница входа.

\textbf{Постусловие:} пользователь успешно авторизован и перенаправлен в личный кабинет либо на главную страницу системы.

\textbf{Основной успешный сценарий:}

\begin{enumerate}
	\item Пользователь открывает страницу входа.
	\item Вводит логин и пароль.
	\item Нажимает кнопку «Войти».
	\item Система выполняет поиск пользователя в базе данных.
	\item Выполняется проверка пароля.
	\item Создаётся пользовательская сессия.
	\item Пользователь получает доступ к функциям системы.
\end{enumerate}

\textbf{Альтернативный сценарий (неверные данные):}

\begin{enumerate}
	\item Пользователь вводит неверный логин или пароль.
	\item Система отображает сообщение об ошибке авторизации.
	\item Пользователь остаётся на странице входа.
\end{enumerate}

\textbf{Альтернативный сценарий (пустые поля):}

\begin{enumerate}
	\item Пользователь не заполняет одно из обязательных полей.
	\item Система выводит сообщение о необходимости заполнения формы.
\end{enumerate}

\subsubsection{Вариант использования «Выход из системы»}

\textbf{Заинтересованные лица и их требования:} авторизованный пользователь желает завершить работу с системой.

\textbf{Предусловие:} пользователь авторизован.

\textbf{Постусловие:} пользовательская сессия завершена.

\textbf{Основной успешный сценарий:}

\begin{enumerate}
	\item Пользователь нажимает кнопку «Выход».
	\item Система удаляет активную сессию.
	\item Пользователь перенаправляется на страницу входа.
\end{enumerate}

\subsubsection{Вариант использования «Загрузка документа на проверку»}

\textbf{Заинтересованные лица и их требования:} пользователь желает отправить документ в формате DOCX для автоматической проверки.

\textbf{Предусловие:} пользователь авторизован.

\textbf{Постусловие:} файл загружен в систему и передан модулю проверки.

\textbf{Основной успешный сценарий:}

\begin{enumerate}
	\item Пользователь открывает страницу загрузки.
	\item Выбирает файл формата DOCX.
	\item Нажимает кнопку «Загрузить».
	\item Система сохраняет файл во временное хранилище.
	\item Запускается процедура проверки документа.
\end{enumerate}

\textbf{Альтернативный сценарий (неподдерживаемый формат):}

\begin{enumerate}
	\item Пользователь выбирает файл иного формата.
	\item Система выводит сообщение об ошибке.
	\item Загрузка не выполняется.
\end{enumerate}

\textbf{Альтернативный сценарий (файл не выбран):}

\begin{enumerate}
	\item Пользователь нажимает кнопку загрузки без выбора файла.
	\item Система предлагает выбрать документ.
\end{enumerate}

\subsubsection{Вариант использования «Автоматическая проверка документа»}

\textbf{Заинтересованные лица и их требования:} пользователь хочет увидеть, где именно в его тексте нарушены стандарты оформления.

\textbf{Предусловие:} файл загружен на сервер и доступен для чтения.

\textbf{Постусловие:} алгоритм сгенерировал итоговую сводку по проверке.

\textbf{Основной успешный сценарий:}

\begin{enumerate}
	\item Парсер системы разбирает внутреннюю структуру файла.
	\item Выполняется проверка параметров оформления.
	\item Модуль фиксирует все отклонения от заданных требований.
	\item Платформа собирает найденные дефекты в единый список.
\end{enumerate}

\textbf{Альтернативный сценарий (ошибка обработки файла):}

\begin{enumerate}
	\item Документ повреждён либо содержит некорректную структуру.
	\item Сервер принудительно останавливает процесс проверки.
	\item В интерфейсе всплывает уведомление о невозможности прочитать документ.
\end{enumerate}

\subsubsection{Вариант использования «Получение отчёта об ошибках»}

\textbf{Заинтересованные лица и их требования:} пользователь желает ознакомиться с найденными нарушениями.

\textbf{Предусловие:} проверка завершена.

\textbf{Постусловие:} пользователю отображён отчёт.

\textbf{Основной успешный сценарий:}

\begin{enumerate}
	\item После завершения проверки система формирует отчёт.
	\item Пользователь открывает страницу результата.
	\item Отображается список ошибок с пояснениями.
\end{enumerate}

\textbf{Альтернативный сценарий (ошибок не обнаружено):}

\begin{enumerate}
	\item Проверка завершена без замечаний.
	\item Система отображает сообщение о корректности документа.
\end{enumerate}

\subsubsection{Вариант использования «Скачивание документа с выделенными ошибками»}

\textbf{Заинтересованные лица и их требования:} пользователь желает скачать исправляемый файл с визуально отмеченными ошибками.

\textbf{Предусловие:} документ содержит нарушения.

\textbf{Постусловие:} пользователю предоставлен файл DOCX для скачивания.

\textbf{Основной успешный сценарий:}

\begin{enumerate}
	\item Пользователь нажимает кнопку скачивания.
	\item Система формирует копию документа.
	\item Ошибки подчёркиваются или выделяются.
	\item Пользователь скачивает файл.
\end{enumerate}

\textbf{Альтернативный сценарий (файл не сформирован):}

\begin{enumerate}
	\item Во время генерации файла возникает ошибка.
	\item Система выводит сообщение о невозможности подготовки документа.
\end{enumerate}

\subsubsection{Вариант использования «Автоматическое сохранение корректного документа»}

\textbf{Заинтересованные лица и их требования:} пользователь желает сохранить корректно оформленную работу в архиве системы.

\textbf{Предусловие:} ошибок в документе не обнаружено.

\textbf{Постусловие:} файл помещён в централизованное хранилище.

\textbf{Основной успешный сценарий:}

\begin{enumerate}
	\item Проверка завершается без ошибок.
	\item Система сохраняет документ в архив.
	\item Информация о работе заносится в базу данных.
\end{enumerate}

\textbf{Альтернативный сценарий (ошибка сохранения):}

\begin{enumerate}
	\item Возникает ошибка доступа к хранилищу либо базе данных.
	\item Система фиксирует ошибку в журнале событий.
	\item Пользователю выводится уведомление.
\end{enumerate}

\subsubsection{Вариант использования «Просмотр архива документов»}

\textbf{Заинтересованные лица и их требования:} пользователю необходим доступ к истории ранее проверенных файлов.

\textbf{Предусловие:} сессия пользователя активна (пройдена авторизация).

\textbf{Постусловие:} на странице отрисована таблица с историей загрузок.

\textbf{Основной успешный сценарий:}

\begin{enumerate}
	\item Пользователь переходит во вкладку архивных записей.
	\item Сервер извлекает из базы данных перечень прошлых проверок.
	\item Пользователь находит интересующий материал и сохраняет его на локальный диск.
\end{enumerate}

\textbf{Альтернативный сценарий (архив пуст):}

\begin{enumerate}
	\item В базе данных нет ни одной записи, привязанной к текущему профилю.
	\item Интерфейс выводит экранную заглушку с текстом «Архив пуст».
\end{enumerate}

\subsubsection{Вариант использования «Формирование титульного листа»}

\textbf{Заинтересованные лица и их требования:} пользователь желает автоматически создать титульный лист по шаблону вуза.

\textbf{Предусловие:} пользователь авторизован.

\textbf{Постусловие:} сформирован титульный лист.

\textbf{Основной успешный сценарий:}

\begin{enumerate}
	\item Пользователь открывает форму генерации титульного листа.
	\item Заполняет необходимые поля.
	\item Нажимает кнопку формирования.
	\item Система создаёт документ DOCX.
\end{enumerate}

\textbf{Альтернативный сценарий (не заполнены обязательные поля):}

\begin{enumerate}
	\item Пользователь оставляет обязательные поля пустыми.
	\item Система сообщает о необходимости заполнения данных.
\end{enumerate}

\subsubsection{Вариант использования «Просмотр статистики»}

\textbf{Заинтересованные лица и их требования:} администратор желает получить статистические сведения о работе системы.

\textbf{Предусловие:} администратор авторизован.

\textbf{Постусловие:} отображён аналитический отчёт.

\textbf{Основной успешный сценарий:}

\begin{enumerate}
	\item Администратор открывает раздел статистики.
	\item Система подсчитывает показатели.
	\item Отображаются данные по количеству проверок, ошибок, сохранённых документов и пользователям.
\end{enumerate}

\textbf{Альтернативный сценарий (недостаточно прав):}

\begin{enumerate}
	\item Обычный пользователь пытается открыть раздел статистики.
	\item Система запрещает доступ.
\end{enumerate}

\subsubsection{Вариант использования «Управление пользователями»}

\textbf{Заинтересованные лица и их требования:} администратор желает управлять учётными записями пользователей.

\textbf{Предусловие:} администратор авторизован.

\textbf{Постусловие:} учётные записи обновлены.

\textbf{Основной успешный сценарий:}

\begin{enumerate}
	\item Администратор открывает раздел пользователей.
	\item Добавляет, редактирует или удаляет записи.
	\item Система сохраняет изменения.
\end{enumerate}

\textbf{Альтернативный сценарий (дублирование логина):}

\begin{enumerate}
	\item При создании пользователя введён уже существующий логин.
	\item Система выводит сообщение об ошибке.
\end{enumerate}

\subsubsection{Вариант использования «Настройка параметров системы»}

\textbf{Заинтересованные лица и их требования:} администратор желает изменить правила проверки документов и параметры работы системы.

\textbf{Предусловие:} учетная запись обладает правами администратора.

\textbf{Постусловие:} новые параметры конфигурации сохранены в систему.

\textbf{Основной успешный сценарий:}

\begin{enumerate}
	\item Администратор переходит в панель управления системой.
	\item Изменяет параметры.
	\item Подтверждает сохранение через пользовательский интерфейс.
	\item Система принимает новые значения.
\end{enumerate}

\textbf{Альтернативный сценарий (некорректные параметры):}

\begin{enumerate}
	\item В формы ввода переданы значения, нарушающие бизнес-логику.
	\item Система отменяет сохранение и выводит сообщение об ошибке
\end{enumerate}

\subsection{Требования к надёжности}

Спроектированная информационно-аналитическая система обязана выдерживать пиковые нагрузки без программных сбоев, а также полностью исключать риск потери загруженных файлов или сгенерированных метрик.

\textbf{Целостность данных.}
Любые изменения состояния платформы — будь то загрузка нового документа, логирование результатов проверки, создание учетной записи или экспорт статистики — должны выполняться строго в рамках атомарных транзакций.При возникновении ошибки операция должна быть отменена без повреждения существующих данных.

\textbf{Устойчивость к сбоям.}
При аварийном завершении работы сервера, отключении электропитания или программном сбое система должна сохранять ранее записанные данные. После повторного запуска система должна продолжить работу в штатном режиме.

\textbf{Контроль корректности входных данных.}
На уровне приложения должна выполняться проверка корректности вводимых пользователями данных: обязательность заполнения полей, допустимые форматы файлов, ограничения длины строковых значений, корректность дат и иных параметров.

\textbf{Логирование ошибок.}
Критические ошибки, исключения и события работы системы должны фиксироваться в журнале событий. Записи журнала должны содержать дату, время, уровень важности и описание ошибки.

\textbf{Резервное копирование.}
База данных и архив документов должны допускать резервное копирование штатными средствами используемой СУБД и файловой системы. Восстановление данных должно выполняться из резервной копии.

\subsection{Требования к информационной безопасности}

Система должна обеспечивать защиту данных от несанкционированного доступа и разграничение прав пользователей.

\textbf{Авторизация пользователей.}
Доступ к функциям системы должен предоставляться только после успешного ввода логина и пароля.

\textbf{Хранение паролей.}
Пароли пользователей не должны храниться в открытом виде. Для хранения необходимо использовать криптографическое хеширование.

\textbf{Разграничение прав доступа.}
В системе должны поддерживаться роли пользователей:

\begin{itemize}
	\item пользователь --- загрузка документов, просмотр результатов, генерация титульных листов;
	\item администратор --- полный доступ, включая управление пользователями и просмотр статистики.
\end{itemize}

\textbf{Проверка доступа на серверной стороне.}
Права пользователя должны проверяться сервером при выполнении каждого защищённого действия.

\textbf{Защита от SQL-инъекций.}
Доступ к базе данных должен осуществляться с использованием ORM либо параметризованных запросов.

\textbf{Защита файлового хранилища.}
Загруженные документы и архивные материалы должны храниться в каталогах, недоступных для прямого неавторизованного обращения.

\subsection{Требования к программной и технической совместимости}

Система должна функционировать в типовой вычислительной среде и не предъявлять завышенных требований к аппаратным ресурсам.

\textbf{Операционная система сервера:}

\begin{itemize}
	\item Windows 10/11.
\end{itemize}

\textbf{Аппаратные требования:}

\begin{itemize}
	\item процессор с тактовой частотой от 2 ГГц;
	\item оперативная память не менее 4 ГБ;
	\item свободное дисковое пространство от 1 ГБ без учёта пользовательских файлов.
\end{itemize}

\textbf{Клиентская часть.}
Для работы пользователей необходимы текстовый редактор Word и современный веб-браузер:

\begin{itemize}
	\item Яндекс Браузер;
	\item Google Chrome;
	\item Mozilla Firefox;
	\item Microsoft Edge.
\end{itemize}



\subsection{Требования к оформлению документации}

Разработка программной документации и программного изделия должна выполняться в соответствии с требованиями действующих нормативных документов ГОСТ 19.102--77 и ГОСТ 34.601--90. Единая система программной документации.